\section{loss 的几种形式}

\subsection{naive loss}

最直接的 loss 是
\[
\mathcal{L}_{\text{naive}} = -\mathbb{E}[\log \pi(a \mid s)\,\delta].
\]
它就是把 log probability 和 reward 信号直接乘起来。


\subsection{old policy ratio}

如果引入旧策略 $\pi_{\text{old}}$，则会出现 ratio：
\[
\frac{\pi(a\mid s)}{\pi_{\text{old}}(a\mid s)}.
\]
脚本里把它写成 log prob 差：
\[
\exp(\log \pi - \log \pi_{\text{old}}).
\]

\begin{warningbox}{旧策略必须 freeze}
`old policy` 只能当常量用，不能让梯度穿过去。否则 ratio 的意义会被破坏，整个优化会偏掉。
\end{warningbox}


\subsection{clipping}

和 PPO 类似，ratio 太大或太小时都不稳定，所以会做 clipping：
\[
r = \frac{\pi}{\pi_{\text{old}}}, \qquad
\tilde r = \mathrm{clip}(r, 1-\epsilon, 1+\epsilon).
\]
最后用 unclipped 和 clipped 的较小者构成保守更新。


\subsection{KL penalty}

如果担心 RL 把模型带偏太远，还可以加 reference model 的 KL penalty：
\[
\mathrm{KL}(p\|q) = \mathbb{E}_{x\sim p}\!\left[\log \frac{p(x)}{q(x)}\right].
\]
脚本里用了一个数值上可算的估计：
\[
\mathbb{E}_p\!\left[\frac{q}{p} - \log\frac{q}{p} - 1\right].
\]

\begin{importantbox}{KL penalty 的作用}
它是“别忘了原来的模型能力”的安全带。你想让模型获得新能力，但又不想把旧能力完全冲掉。
\end{importantbox}


